% GENERATED FILE. Edit canonical lessons, then run npm run generate:books.
% canonical-source-hash: fnv1a64:cc4f9bd5b46e6d5f
% canonical-lessons: ZH-C490-guanli, ZH-C490-hezuo, ZH-C490-yeban, ZH-C490-xinku, ZH-C490-kaichu

\chapter{Pressure at Work}
\label{ch:zh-pressure-at-work}
\hlchaptermodality{490}

\begin{chapteropening}
\textbf{By the end of this chapter:} \emph{I can talk about managing, working together, night shifts, hard work and being fired.}

Complete the last lesson of chapter 490: I can talk about managing, working together, night shifts, hard work and being fired.
\end{chapteropening}

\section[guǎnlǐ]{guǎnlǐ — to manage; to be in charge of}

\label{lesson:ZH-C490-guanli}



\begin{quote}
Before the new one: say the Mandarin for a position in a firm, then the Mandarin for to get promoted.
\end{quote}

\subsection*{You'll want to know: guǎnlǐ}
\emph{guǎnlǐ} — \textquotedblleft{}to manage; to be in charge of\textquotedblright{}.

\begin{quote}
\emph{tā guǎnlǐ yì jiā fàndiàn} — She manages a restaurant.
\end{quote}

Its characters come later; for now it is a word for the ear.

Two third tones meet in this word: guǎn is said with a rising tone before lǐ, and the spelling keeps both marks.

\begin{cousinweb}
\emph{guǎnlǐ} is built from \emph{guǎn} 'to be in charge' and \zh{理} 'to put in order', both inherited from Middle Chinese, \emph{kwanX} and \emph{liX} (Guangyun, in Baxter's transcription), which glosses the first 'a pipe (a wind instrument)' and 'to be in charge', and the second 'to work jade' and 'a pattern'.
\end{cousinweb}

\begin{grammarlens}[title={What you've built}]
One more word for saying how it was.
\end{grammarlens}

\subsection*{Your turn}
\begin{itemize}
\raggedright
  \item \emph{Say it:} \emph{guǎnlǐ}
  \item \emph{Say it:} \emph{guǎnlǐ}, once more
  \item \emph{Say it:} \emph{tā guǎnlǐ yì jiā fàndiàn} — She manages a restaurant.
  \item \emph{Recall:} say the Mandarin for a position in a firm, then the Mandarin for to get promoted, then say \emph{guǎnlǐ} again
\end{itemize}

\begin{tcolorbox}[breakable,colback=teal!4,colframe=teal!35!black,title={Before you move on}]
What does guǎnlǐ mean? (\textquotedblleft{}To manage; to be in charge of\textquotedblright{}.) Say it once more.
\end{tcolorbox}
\section[hézuò]{hézuò — to work together; to cooperate}

\label{lesson:ZH-C490-hezuo}



\begin{quote}
Before the new one: say the Mandarin for to get promoted, then the Mandarin for to manage.
\end{quote}

\subsection*{You'll want to know: hézuò}
\emph{hézuò} — \textquotedblleft{}to work together; to cooperate\textquotedblright{}.

\begin{quote}
\emph{gēn tā hézuò wǒ gǎnjué hěn kāixīn} — I really enjoy working with him.
\end{quote}

Its characters come later; for now it is a word for the ear.

\begin{cousinweb}
\emph{hézuò} is built from \emph{hé} 'to join' and \zh{作} 'to do', both inherited from Middle Chinese, \emph{hop} and \emph{tsak} (Guangyun, in Baxter's transcription), which glosses the first 'to join together' and the second 'to do' and 'to make'.
\end{cousinweb}

\begin{grammarlens}[title={What you've built}]
One more word for saying how it was.
\end{grammarlens}

\subsection*{Your turn}
\begin{itemize}
\raggedright
  \item \emph{Say it:} \emph{hézuò}
  \item \emph{Say it:} \emph{hézuò}, once more
  \item \emph{Say it:} \emph{gēn tā hézuò wǒ gǎnjué hěn kāixīn} — I really enjoy working with him.
  \item \emph{Recall:} say the Mandarin for to get promoted, then the Mandarin for to manage, then say \emph{hézuò} again
\end{itemize}

\begin{tcolorbox}[breakable,colback=teal!4,colframe=teal!35!black,title={Before you move on}]
What does hézuò mean? (\textquotedblleft{}To work together; to cooperate\textquotedblright{}.) Say it once more.
\end{tcolorbox}
\section[yèbān]{yèbān — night shift}

\label{lesson:ZH-C490-yeban}



\begin{quote}
Before the new one: say the Mandarin for to manage, then the Mandarin for to work together.
\end{quote}

\subsection*{You'll want to know: yèbān}
\emph{yèbān} — \textquotedblleft{}night shift\textquotedblright{}.

\begin{quote}
\emph{zuò le yì nián yèbān yǐhòu tā chángcháng shīmián} — After a year of night shifts, he often couldn't sleep.
\end{quote}

Its characters come later; for now it is a word for the ear.

\begin{cousinweb}
\emph{yèbān} is built from \emph{yè} 'night' and \emph{bān} 'a shift', both inherited from Middle Chinese, \emph{yaeH} and \emph{paen} (Guangyun, in Baxter's transcription), which glosses the first 'dusk' and the second 'to divide jade tokens'; the sense 'a turn of work' is not given there.
\end{cousinweb}

\begin{grammarlens}[title={What you've built}]
One more word for saying how it was.
\end{grammarlens}

\subsection*{Your turn}
\begin{itemize}
\raggedright
  \item \emph{Say it:} \emph{yèbān}
  \item \emph{Say it:} \emph{yèbān}, once more
  \item \emph{Say it:} \emph{zuò le yì nián yèbān yǐhòu tā chángcháng shīmián} — After a year of night shifts, he often couldn't sleep.
  \item \emph{Recall:} say the Mandarin for to manage, then the Mandarin for to work together, then say \emph{yèbān} again
\end{itemize}

\begin{tcolorbox}[breakable,colback=teal!4,colframe=teal!35!black,title={Before you move on}]
What does yèbān mean? (\textquotedblleft{}Night shift\textquotedblright{}.) Say it once more.
\end{tcolorbox}
\section[xīnkǔ]{xīnkǔ — exhausting; toilsome (of work)}

\label{lesson:ZH-C490-xinku}



\begin{quote}
Before the new one: say the Mandarin for to work together, then the Mandarin for the night shift.
\end{quote}

\subsection*{You'll want to know: xīnkǔ}
\emph{xīnkǔ} — \textquotedblleft{}exhausting; toilsome (of work)\textquotedblright{}.

\begin{quote}
\emph{gōngzuò yālì dà de shíhou wǒ gǎnjué hěn xīnkǔ} — When the pressure at work is high, I find it really hard going.
\end{quote}

Its characters come later; for now it is a word for the ear.

\begin{cousinweb}
\emph{xīnkǔ} is built from \emph{xīn} 'pungent' and \emph{kǔ} 'bitter', both inherited from Middle Chinese, \emph{sin} and \emph{khuX} (Guangyun, in Baxter's transcription), which glosses the first as the taste of ginger and the second 'coarse' and 'toil'.
\end{cousinweb}

\begin{grammarlens}[title={What you've built}]
One more word for saying how it was.
\end{grammarlens}

\subsection*{Your turn}
\begin{itemize}
\raggedright
  \item \emph{Say it:} \emph{xīnkǔ}
  \item \emph{Say it:} \emph{xīnkǔ}, once more
  \item \emph{Say it:} \emph{gōngzuò yālì dà de shíhou wǒ gǎnjué hěn xīnkǔ} — When the pressure at work is high, I find it really hard going.
  \item \emph{Recall:} say the Mandarin for to work together, then the Mandarin for the night shift, then say \emph{xīnkǔ} again
\end{itemize}

\begin{tcolorbox}[breakable,colback=teal!4,colframe=teal!35!black,title={Before you move on}]
What does xīnkǔ mean? (\textquotedblleft{}Exhausting; toilsome (of work)\textquotedblright{}.) Say it once more.
\end{tcolorbox}
\section[kāichú]{kāichú — to fire; to dismiss (from a job or school)}

\label{lesson:ZH-C490-kaichu}



\begin{quote}
Before the new one: say the Mandarin for the night shift, then the Mandarin for exhausting (of work).
\end{quote}

\subsection*{You'll want to know: kāichú}
\emph{kāichú} — \textquotedblleft{}to fire; to dismiss (from a job or school)\textquotedblright{}.

\begin{quote}
\emph{tā chángcháng chídào suǒyǐ lǎobǎn kāichú le tā} — He was often late, so the boss fired him.
\end{quote}

Its characters come later; for now it is a word for the ear.

\begin{cousinweb}
\emph{kāichú} is built from \zh{开} 'to open' and \emph{chú} 'to remove', both inherited from Middle Chinese, \emph{khoj} and \emph{drjo} (Guangyun, in Baxter's transcription), which glosses the first 'to release' and the second 'to take away'.
\end{cousinweb}

\begin{grammarlens}[title={What you've built}]
One more word for saying how it was.
\end{grammarlens}

\subsection*{Your turn}
\begin{itemize}
\raggedright
  \item \emph{Say it:} \emph{kāichú}
  \item \emph{Say it:} \emph{kāichú}, once more
  \item \emph{Say it:} \emph{tā chángcháng chídào suǒyǐ lǎobǎn kāichú le tā} — He was often late, so the boss fired him.
  \item \emph{Recall:} say the Mandarin for the night shift, then the Mandarin for exhausting (of work), then say \emph{kāichú} again
\end{itemize}

\begin{tcolorbox}[breakable,colback=teal!4,colframe=teal!35!black,title={Before you move on}]
What does kāichú mean? (\textquotedblleft{}To fire; to dismiss (from a job or school)\textquotedblright{}.) Say it once more.
\end{tcolorbox}
